%%%%%%%%%%%%%%%%%%%%%%%%%%%%%%%%%%%%%%%%%%%%%%%%%%%%%%%%%%%%%%%%%%%%%%%%%%%
% Handout und Uebungsblatt (Teil A, B, C) zu Lehreinheit 4 (Quellenkritik)
% Im Skript am Kapitelende, in handouts.tex als Druckvorlage.
%%%%%%%%%%%%%%%%%%%%%%%%%%%%%%%%%%%%%%%%%%%%%%%%%%%%%%%%%%%%%%%%%%%%%%%%%%%
\begin{handout}{Quellenkritik}

\textbf{Worum es heute ging.} In reinen Zufallspunkten haben Sie Muster
gesehen, und ein erfundener, gut gestalteter Artikel wirkte
glaubwürdiger als eine echte, unscheinbare Studie. Das Gehirn ergänzt
Muster und Glaubwürdigkeit aus Oberflächenmerkmalen. Quellenkritik
setzt an die Stelle des Eindrucks eine Prüfung.

\textbf{Kernbegriffe}
\begin{description}
  \item[Apophänie] die Neigung, in Zufall Muster zu sehen
    \parencite{conrad1958schizophrenie}.
  \item[Bestätigungsfehler] Belege, die uns recht geben, werden weniger
    streng geprüft \parencite{nickerson1998ubiquitous}.
  \item[Primärstudie / Review / Meta-Analyse] eigene Daten /
    Zusammenfassung mehrerer Studien / statistische Zusammenfassung ihrer
    Effektgrößen. \textbf{Graue Literatur}: Preprints, Berichte,
    Abschlussarbeiten.
  \item[Impact Factor, h-Index] zählen Zitate, also Aufmerksamkeit, nicht
    Richtigkeit \parencite{seglen1997impact,hirsch2005index}.
  \item[Laterales Lesen] die Quelle früh verlassen und nachschlagen, was
    andere über sie sagen \parencite{wineburg2019lateral}.
  \item[Prebunking] Manipulationstechniken vorab in schwacher Dosis
    kennenlernen \parencite{roozenbeek2022inoculation}.
\end{description}

\textbf{Fünf Kriterien und SIFT}

\begin{center}
\small
\begin{tabular}{@{}p{0.44\textwidth}p{0.48\textwidth}@{}}
\toprule
\textbf{Kriterien der Quellenkritik} & \textbf{SIFT} \parencite{caulfield2023verified} \\
\midrule
Herkunft: Wer spricht, mit welcher Expertise? & \textbf{S}top: innehalten, Ziel klären \\
Intention: Wozu, wer bezahlt? & \textbf{I}nvestigate: außerhalb nachlesen, wer die Quelle ist \\
Methode: Wie gewonnen? & \textbf{F}ind better coverage: bessere Berichte suchen \\
Peer Review: Wer hat geprüft? & \textbf{T}race claims: zum Original zurückverfolgen \\
Aktualität: widerrufen, überholt, repliziert? & \\
\bottomrule
\end{tabular}
\end{center}

\begin{merkbox}
Gutes Layout, eine präzise Zahl und häufige Wiederholung machen eine
Aussage glaubwürdiger, aber nicht wahrer. Fragen Sie bei jeder Zahl: Wer
hat das gezählt, und wo steht das?
\end{merkbox}

\textbf{Quellenkritik in fünf Schritten:} (1)~innehalten, (2)~Herkunft
und Absicht klären, (3)~Art und Methode bestimmen, (4)~Spur zum Original
verfolgen, (5)~Gegenprobe: unabhängige Bestätigung, Widerruf,
Gegenposition.

\textbf{Zum Weiterlesen:} \textcite{wineburg2019lateral};
\textcite{caulfield2023verified}; \textcite{seglen1997impact}.
\end{handout}

\begin{uebungsblatt}{Quellenkritik und SIFT}

\aufgabe{SIFT-Raster für drei Quellen}[in der Stunde, Kleingruppe, 30 Min.; Fertigstellung als Hausaufgabe]
Die Lehrperson gibt drei Quellen vor (Übungsblatt~C): (1)~eine
Peer-Review-Studie, (2)~einen Zeitungsartikel, (3)~einen Wikipedia-Eintrag. Tragen Sie
zuerst die genauen Angaben ein, dann Ihre Befunde zu jedem Schritt.
Verfolgen Sie in Schritt~T mindestens eine Zahl bis zur Originalstudie.
Zur Art der Quelle: \emph{primär} = eigene Daten (Studie),
\emph{sekundär} = fasst Studien zusammen (Review, Meta-Analyse),
\emph{tertiär} = berichtet darüber (Presse, Lexikon), \emph{grau} = ohne
Verlag oder Begutachtung veröffentlicht (Preprint, Bericht,
Abschlussarbeit).

\begin{center}
\small
\renewcommand{\arraystretch}{1.3}
\begin{tabular}{|>{\raggedright\arraybackslash}p{0.15\textwidth}|p{0.24\textwidth}|p{0.24\textwidth}|p{0.24\textwidth}|}
\hline
 & \textbf{(1) Studie} & \textbf{(2) Zeitungsartikel} & \textbf{(3) Wikipedia} \\
\hline
Angaben (Autor, Titel, Datum, Link) & \rule{0pt}{14mm} & & \\ \hline
\textbf{S}top: Was will ich wissen? Erster Eindruck? & \rule{0pt}{16mm} & & \\ \hline
\textbf{I}nvestigate: Wer ist die Quelle? Was sagen andere? & \rule{0pt}{18mm} & & \\ \hline
\textbf{F}ind: Bessere oder unabhängige Berichte? & \rule{0pt}{16mm} & & \\ \hline
\textbf{T}race: Woher stammt die zentrale Zahl/Aussage? & \rule{0pt}{18mm} & & \\ \hline
Art der Quelle (primär, sekundär, tertiär, grau) & \rule{0pt}{9mm} & & \\ \hline
Gewicht (hoch/mittel/gering) und Begründung & \rule{0pt}{16mm} & & \\ \hline
\end{tabular}
\end{center}
\end{uebungsblatt}

\begin{uebungsblattb}{2}{Schöne Fälschung und Quellenkritik}

\aufgabe{Anatomie einer schönen Fälschung}[in der Stunde nach der Auflösung, Einzelarbeit, 10 Min.]
Die folgende Meldung ist \textbf{fiktiv} und für diese Übung konstruiert,
nach dem Muster der \enquote{schönen Fälschung} aus dem Staunmoment.
Markieren Sie alle Merkmale, die Glaubwürdigkeit \emph{erzeugen}, und
notieren Sie zu jedem, ob es etwas \emph{belegt}.

\begin{center}
\fbox{\parbox{0.9\textwidth}{\small
\textbf{Jonglieren macht schlau: Zehn Minuten am Morgen steigern die
Merkfähigkeit um 23,7 Prozent}\\[1mm]
Wer den Tag mit Jonglierbällen beginnt, lernt schneller. Das zeigt eine
neue Studie, über die Fachleute seit Wochen diskutieren. \enquote{Die
Effekte sind beeindruckend}, sagt Prof.\ Dr.\ H.\ Lindner, Lernforscher und
Autor mehrerer Ratgeber. Die Teilnehmenden, die vor dem Vokabellernen
jonglierten, behielten im Schnitt 23,7\,\% mehr Wörter. Immer mehr Schulen
denken bereits über Jonglierpausen nach. Das passt zu dem, was Pädagogen
schon lange vermuten: Bewegung und Denken gehören zusammen.\\[1mm]
{\footnotesize\itshape (Fiktives Lehrmaterial, LE~4. Person, Studie und Zahlen sind erfunden.)}}}
\end{center}

\begin{center}
\renewcommand{\arraystretch}{1.9}
\begin{tabular}{|p{0.42\textwidth}|p{0.48\textwidth}|}
\hline
\textbf{Merkmal} & \textbf{Belegt es etwas? Was fehlt?} \\
\hline
 & \\ \hline
 & \\ \hline
 & \\ \hline
 & \\ \hline
\end{tabular}
\end{center}

\aufgabe{Quellenkritik in fünf Schritten}[Hausaufgabe]
Tragen Sie in der Sicherung die fünf Schritte, die die Gruppe an der
Tafel erarbeitet, in die linke Spalte ein (zum Abgleich: Handout). Wenden
Sie die Checkliste dann auf zwei Einträge Ihrer annotierten Bibliographie
aus LE~3 an. Haken Sie ab, was Sie geprüft haben, und notieren Sie, ob
sich Ihr Urteil geändert hat.

\begin{center}
\renewcommand{\arraystretch}{1.6}
\begin{tabular}{|p{0.5\textwidth}|c|c|}
\hline
\textbf{Schritt} & \textbf{Quelle A} & \textbf{Quelle B} \\
\hline
1 \rule{0pt}{7mm} & $\square$ & $\square$ \\ \hline
2 \rule{0pt}{7mm} & $\square$ & $\square$ \\ \hline
3 \rule{0pt}{7mm} & $\square$ & $\square$ \\ \hline
4 \rule{0pt}{7mm} & $\square$ & $\square$ \\ \hline
5 \rule{0pt}{7mm} & $\square$ & $\square$ \\ \hline
\end{tabular}
\end{center}
Hat sich Ihr Urteil geändert? Warum?
\schreiblinien{3}

\aufgabe{Diskussion: 500 Zitierungen}[Hausaufgabe, max.\ 1 Seite]
\enquote{Eine Studie hat 500 Zitierungen. Was sagt das über ihre
Qualität?} Schreiben Sie eine strukturierte Gegenargumentation gegen die
Annahme, viele Zitate bedeuteten hohe Qualität. Gliederung: (a)~die
These in ihrer stärksten Form, (b)~mindestens drei Gegenargumente mit je
einem Beleg oder Beispiel, (c)~was Zitate \emph{doch} aussagen, (d)~Ihr
Fazit: Wie würden Sie die Qualität stattdessen prüfen?

\aufgabe{Tagebuch des Staunens}[Hausaufgabe]
Beschreiben Sie ein Muster, das Ihnen diese Woche aufgefallen ist
(Zufall, Statistik, Nachrichten, Alltag). Was haben Sie erwartet, was
haben Sie gesehen, und wie könnten Sie prüfen, ob das Muster echt ist?

\begin{lehrhinweis}[Erwartungshorizont]
Aufgabe~1: siehe Lehrhinweis in der Stunde; gute
Lösungen nennen in Schritt~I eine \emph{externe} Quelle und in
Schritt~T die Stelle der Studie, aus der die Zahl stammt.
Aufgabe~2: Glaubwürdigkeitsmerkmale sind die präzise Zahl (23,7\,\%),
der Professorentitel, das Zitat, \enquote{eine neue Studie} ohne Angabe,
\enquote{Fachleute diskutieren}, \enquote{immer mehr Schulen}, die
Anknüpfung an Bekanntes (\enquote{was Pädagogen schon lange vermuten}).
Keines belegt etwas: Es fehlen Studie, Stichprobe, Kontrollgruppe,
Messverfahren, Fachgebiet des Experten (Ratgeberautor:
Interessenkonflikt).
Aufgabe~3: In die linke Spalte gehören die fünf Schritte der Sicherung
(innehalten; Herkunft und Absicht; Art und Methode; Spur zum Original;
Gegenprobe), wie im Handout.
Aufgabe~4: Erwartete Gegenargumente: Zitate messen Aufmerksamkeit, nicht
Richtigkeit; auch Kritik und Widerlegung zählen als Zitat; Matthäus-Effekt
\parencite{merton1968matthew}; Fach- und Alterseffekte (ältere Arbeiten
und große Fächer sammeln mehr Zitate); Selbstzitate; Zitate werden oft
ungelesen übernommen. Zugestanden werden sollte: Eine viel zitierte
Arbeit ist in der Fachdiskussion präsent und meist ein guter Einstieg.
Stärkere Lösungen schlagen Prüfwege vor: Methode lesen, Replikationen
und Widerrufe suchen, in welchem Sinn zitiert wird
\parencite{dora2012,hicks2015leiden}.
\end{lehrhinweis}
\end{uebungsblattb}

\begin{uebungsblattc}{Drei Quellen für die SIFT-Übung}

Diese drei Quellen bearbeiten Sie mit dem SIFT-Raster (Übungsblatt~A,
Aufgabe~1). Sie hängen zusammen: Zeitungsartikel und Wikipedia-Eintrag
berufen sich auf dieselbe Studie. Öffnen Sie die Links über die QR-Codes
und tragen Sie die Angaben selbst ins Raster ein.

\newcommand{\LEvierQuelle}[4]{%
\par\medskip\noindent
\begin{minipage}[t]{0.76\textwidth}
{\bfseries\color{rwucyandark}#1}\par\smallskip
#2\par\smallskip
{\small\url{#3}}
\end{minipage}\hfill
\begin{minipage}[t]{0.2\textwidth}\vspace{0pt}\raggedleft
\qrcode[nolink,height=2.6cm]{#4}
\end{minipage}\par\medskip}

\LEvierQuelle{(1) Die Studie}
{White, M.\,P., Alcock, I., Grellier, J., Wheeler, B.\,W., Hartig, T.,
Warber, S.\,L., Bone, A., Depledge, M.\,H. \& Fleming, L.\,E. (2019).
Spending at least 120 minutes a week in nature is associated with good
health and wellbeing. \emph{Scientific Reports, 9}, Artikel~7730. Frei
zugänglich (Open Access).}
{https://doi.org/10.1038/s41598-019-44097-3}
{https://doi.org/10.1038/s41598-019-44097-3}

\LEvierQuelle{(2) Der Zeitungsartikel}
{DER SPIEGEL (13.\,6.\,2019). Waldbaden und Spazierengehen: Ab zwei
Stunden pro Woche macht die Natur gesund, sagen Forscher. Ressort
Gesundheit.}
{https://www.spiegel.de/gesundheit/ernaehrung/120-minuten-natur-pro-woche-sind-gesund-a-1272293.html}
{https://www.spiegel.de/gesundheit/ernaehrung/120-minuten-natur-pro-woche-sind-gesund-a-1272293.html}

\LEvierQuelle{(3) Der Wikipedia-Eintrag}
{\enquote{Health}, englischsprachige Wikipedia, Fassung vom 26.\,9.\,2026
(permanenter Link, damit alle Gruppen dieselbe Fassung prüfen). Suchen Sie
den Satz, der die Studie zitiert (Abschnitt zur Umwelt als
Gesundheitsfaktor).}
{https://en.wikipedia.org/w/index.php?title=Health&oldid=1376907986}
{https://en.wikipedia.org/w/index.php?title=Health\&oldid=1376907986}

\LEvierQuelle{Zusatz für Schritt T: die Pressemitteilung}
{University of Exeter (Juni 2019). Two hours a week is key dose of nature
for health and wellbeing. Pressemitteilung der Hochschule, an der die
Studie entstand.}
{https://news-archive.exeter.ac.uk/2019/june/title_720059_en.html}
{https://news-archive.exeter.ac.uk/2019/june/title_720059_en.html}

\textbf{Für schnelle Gruppen:} Gibt es in der deutschsprachigen Wikipedia
einen Eintrag, der diese Studie zitiert? Was folgt daraus für Ihre Suche
in Schritt~F?

\begin{lehrhinweis}[Erwartungshorizont zu den drei Quellen]
Die Kette ist geprüft (Oktober 2026): Der SPIEGEL nennt die University of
Exeter und \emph{Scientific Reports}, verlinkt die Studie aber nicht; die
Wikipedia zitiert sie mit DOI. \textbf{S:} Die Überschrift ist kausal und
nennt eine genaue Dosis (\enquote{ab zwei Stunden}); das klingt gemessen
und passt zu Bekanntem (150~Minuten Bewegung, \enquote{5 am Tag}), also
Vorsicht vor dem Bestätigungsfehler. \textbf{I:} Der Artikel ist nicht
namentlich gezeichnet und folgt eng der Pressemitteilung; die Hochschule
hat ein Interesse an Aufmerksamkeit. \textbf{F:} Im SPIEGEL-Text selbst
steht es vorsichtiger als in der Überschrift (\enquote{relativ klein},
umgekehrte Wirkrichtung nicht ausgeschlossen). Die deutschsprachige
Wikipedia zitiert die Studie nicht; der englische Eintrag
\enquote{Health} schreibt \enquote{It has been demonstrated} und
\enquote{positive health benefits}, also stärker als die Studie.
\textbf{T:} Die Studie schreibt \enquote{associated}, nicht
\enquote{macht gesund}: Querschnittsbefragung, Selbstauskunft,
$n = 19\,806$. Die Autoren können nicht ausschließen, dass gesündere
Menschen mehr Zeit in der Natur verbringen; die Zeit in der Natur erklärt
nur rund 1\,\% der Varianz. Die Zahl 120 ist zum Teil ein Artefakt: Die
Zeit wurde in 60-Minuten-Blöcke eingeteilt, und 78\,\% der Angaben im
Block 120--179 Minuten lauten genau 120 (runde Stunden). Brücke zur
Apophänie: Eine \enquote{Schwelle} wird in kategorisierte Daten
hineingelesen. \textbf{Gewicht:} Studie mittel (groß, begutachtet, aber
korrelativ, von den Autoren selbst als vorläufig bezeichnet), Zeitung
gering bis mittel, Wikipedia gering (tertiär, überzeichnet).
\end{lehrhinweis}
\end{uebungsblattc}
